\section{Periodic codes and the Correspondence Principle}
\label{sec:correspondence}

The algebraic heart of the Correspondence Principle is a fixed-point
identity for a repeated word.  Its dynamical heart is branch correctness.
Those are independent steps: the first takes place in an affine group, while
the second identifies a formal composition with an iterate of a piecewise
map.

\subsection{Purely periodic digit strings}

Let $\mathbf j=(j_0,\ldots,j_{L-1})$ be a word of length $L\geq1$, and put
\[
 n=\Dig(\mathbf j)=\sum_{a=0}^{L-1}j_ap^a.
\]
Its infinite repetition represents the base-$p$ digit-ring element
\begin{equation}
\label{eq:Bp}
 B_p(\mathbf j)=B_{p,L}(n)
 =n(1+p^L+p^{2L}+\cdots)=\frac{n}{1-p^L}.
\end{equation}
The length $L$ is part of this expression.  When $\mathbf j$ is the shortest
word for $n>0$, write $\lambda_p(n)=L$ and
$B_p(n)=n/(1-p^{\lambda_p(n)})$.

\begin{theorem}[repeated-code fixed-point identity]
\label{thm:repeated-code-fixed-point}
Let $H$ be centered, let $\ell$ be a place of $K$, and set
$M=M_H(\mathbf j)$ and $X=X_H(\mathbf j)$.  Suppose
$\|M\|_\ell<1$.  Then $I-M$ is invertible on $K_\ell$, the rising numen
converges at $B_p(\mathbf j)$, and
\begin{equation}
\label{eq:repeated-code-value}
 X_{H,\ell}(B_p(\mathbf j))=(I-M)^{-1}X.
\end{equation}
This value is the unique fixed point of $H_{\mathbf j}$ in $K_\ell$.
\end{theorem}

\begin{proof}
The $mL$-digit truncation of $B_p(\mathbf j)$ is represented by
$\mathbf j^{\wedge m}$.  By \cref{prop:repeated-word},
\[
 X_H(\mathbf j^{\wedge m})
 =(I+M+\cdots+M^{m-1})X.
\]
The operator series converges in the complete endomorphism algebra because
$\|M\|_\ell<1$, and its sum is $(I-M)^{-1}$.  Truncations ending inside the
next block have the same limit because the remaining finite prefix is
multiplied by $M^m$.  This proves the rising limit.  The fixed-point assertion
is \cref{eq:word-fixed-point}.
\end{proof}

\begin{remark}
In a scalar one-dimensional setting \cref{eq:repeated-code-value} is often
written $X/(1-M)$.  For a general Hydra, $M$ is an endomorphism and the typed
formula is $(I-M)^{-1}X$.
\end{remark}

\begin{proposition}[rational points are the eventually periodic digit strings]
\label{prop:rational-eventually-periodic}
Under the diagonal embedding into the completions at the primes dividing
$p$,
\[
 \Q\cap\Z_p^{\mathrm{dig}}
 =\Z_{(p)}:=\left\{\frac ab\in\Q:\gcd(b,p)=1
                   \text{ after reducing }a/b\right\}.
\]
An element of $\Z_p^{\mathrm{dig}}$ lies in this subring if and only if its
base-$p$ digit sequence is eventually periodic.
\end{proposition}

\begin{proof}
Write $p=\prod_{\substack{\ell\mid p\\ \ell\ {\mathrm{prime}}}}\ell^{e_\ell}$.
The compatible Chinese-remainder
isomorphisms give
\[
 \Z_p^{\mathrm{dig}}
 =\varprojlim_N\Z/p^N\Z
 \xrightarrow{\sim}
 \prod_{\ell\mid p}\varprojlim_N\Z/\ell^{e_\ell N}\Z
 \xrightarrow{\sim}\prod_{\ell\mid p}\Z_\ell.
\]
The subsequence $(e_\ell N)_{N\geq1}$ is cofinal, which proves the second
isomorphism.  Under the resulting diagonal embedding of $\Q$ into
$\prod_{\ell\mid p}\Q_\ell$, a reduced rational belongs to the displayed
product of integer rings exactly when no prime $\ell\mid p$ divides its
denominator.  Hence
$\Q\cap\Z_p^{\mathrm{dig}}=\Z_{(p)}$ with the meaning stated above.

If a digit string has a prefix of length $R$ with value $u$ and thereafter
repeats a block of length $L$ with value
$v=\sum_{h=0}^{L-1}e_hp^h$, then its value in the inverse limit is
\[
 u+p^R v(1+p^L+p^{2L}+\cdots)
 =u+\frac{p^Rv}{1-p^L}.
\]
Because $1-p^L\equiv1\pmod p$, its denominator is coprime to $p$.

Conversely, let $x=a/b$ with $a,b\in\Z$ and $\gcd(b,p)=1$.  For every
$N\geq0$ there is a unique compatible residue
$A_N\in\{0,\ldots,p^N-1\}$ satisfying
$bA_N\equiv a\pmod{p^N}$; these residues define the digit expansion of $x$.
Here $A_0=0$.  Put
\[
 u_N=\frac{a-bA_N}{p^N}\in\Z.
\]
If $d_N$ is the next digit, then
\begin{equation}
\label{eq:rational-digit-state}
 u_{N+1}=\frac{u_N-bd_N}{p},
 \qquad
 d_N\in\mathcal A_p,\qquad bd_N\equiv u_N\pmod p.
\end{equation}
The congruence determines $d_N$ uniquely from $u_N$ because $b$ is a unit
modulo $p$.  Iterating
$|u_{N+1}|\leq(|u_N|+|b|(p-1))/p$ gives
\[
 |u_N|\leq p^{-N}|u_0|+|b|(1-p^{-N})
 \leq |u_0|+|b|.
\]
Thus the integer sequence $(u_N)$ takes values in a finite set, and the
deterministic transition
\cref{eq:rational-digit-state} must eventually repeat.  Its associated
digits are therefore eventually periodic.
\end{proof}

Let $H$ be centered, fix a place $\ell$ at which all branches are bounded,
and suppose $\|M_H(\mathbf j)\|_\ell<1$.  If $\mathfrak z$ has a finite
prefix $\mathbf i$ followed by the infinite repetition of $\mathbf j$, then
\begin{equation}
\label{eq:eventually-periodic-numen}
 X_{H,\ell}(\mathfrak z)
 =H_{\mathbf i}\bigl((I-M_H(\mathbf j))^{-1}X_H(\mathbf j)\bigr).
\end{equation}
This is an affine image of a formal periodic fixed point.  It becomes a
statement about the actual Hydra only after correctness is proved.

\subsection{Classical \texorpdfstring{$2$}{2}-adic conjugacy and the exact
periodic overlap}

The parity-vector conjugacy predates the numen.  Its digit order is
chronological, whereas Siegel's displayed affine word acts from right to
left.  The difference is not cosmetic: it forces a word reversal in the
precise locus where the two constructions have a common rational value.

Let $q$ be any odd integer and extend $T_q$ to $\Z_2$ using the parity
branches in \cref{eq:Tq-branches}.  Let
\[
 S=\vartheta_2:\Z_2\longrightarrow\Z_2,
 \qquad S(z)=\frac{z-[z]_2}{2},
\]
be the binary digit shift.
For every prime $\ell$, write
$\iota_\ell:\Q\hookrightarrow\Q_\ell$ for the canonical field embedding.

\begin{proposition}[the Bernstein--Lagarias conjugacy, with maps typed]
\label{prop:BL-conjugacy-typed}
Define
\begin{equation}
\label{eq:BL-parity-encoder}
 Q_q(x)=\sum_{i=0}^{\infty}\bigl[T_q^{\circ i}(x)\bigr]_2\,2^i
 \qquad(x\in\Z_2).
\end{equation}
If $z\in\Z_2\setminus\{0\}$, let either $I=\{1,\ldots,r\}$ for some $r\geq1$ or
$I=\Z_{\geq1}$, according as the set of one positions is finite or infinite,
and write
\[
 z=\sum_{\ell\in I}2^{d_\ell},
 \qquad d_1\geq0,
 \qquad d_\ell<d_{\ell+1}\quad(\ell,\ell+1\in I).
\]
Define
\begin{equation}
\label{eq:BL-inverse-formula}
 \Phi_q(z)=-\sum_{\ell\in I}q^{-\ell}2^{d_\ell},
 \qquad \Phi_q(0)=0.
\end{equation}
Then $Q_q,\Phi_q:\Z_2\to\Z_2$ are mutually inverse $2$-adic
isometries, hence solenoidal homeomorphisms, and
\begin{equation}
\label{eq:BL-conjugacy-oriented}
 Q_q\circ T_q=S\circ Q_q,
 \qquad
 T_q=\Phi_q\circ S\circ\Phi_q^{-1}.
\end{equation}
Thus $Q_q$ is the chronological parity encoder and
$\Phi_q=Q_q^{-1}$ is the conjugacy.  This is the odd-$(a,b)=(q,1)$
specialization of Bernstein--Lagarias, not a prime-$q$ statement
\cite[Eqs.~(1.3), (1.5), (1.6), (4.1), and~(4.2)]{BernsteinLagarias}.
\end{proposition}

\begin{proof}
For odd $q$, $H_0$ maps $2\Z_2$ to $\Z_2$ and $H_1$ maps
$1+2\Z_2$ to $\Z_2$.  Thus the branch selected by the current parity is
defined at every iterate, every digit in \cref{eq:BL-parity-encoder} lies in
$\{0,1\}$, and the displayed binary series converges.  If
$x\equiv y\pmod{2^N}$, induction through their common parity
branches gives
\[
 T_q^{\circ i}(x)\equiv T_q^{\circ i}(y)
 \pmod{2^{N-i}}
 \qquad(0\leq i\leq N),
\]
so the first $N$ digits of $Q_q(x)$ and $Q_q(y)$ agree.  Conversely, if
those $N$ digits agree and contain $m$ ones, the same-branch affine
calculation gives
\[
 T_q^{\circ N}(x)-T_q^{\circ N}(y)
 =\frac{q^m}{2^N}(x-y).
\]
The left side lies in $\Z_2$ and $q$ is a $2$-adic unit, hence
$x\equiv y\pmod{2^N}$.  Therefore
\begin{equation}
\label{eq:BL-isometry-congruences}
 x\equiv y\pmod{2^N}
 \quad\Longleftrightarrow\quad
 Q_q(x)\equiv Q_q(y)\pmod{2^N},
\end{equation}
which proves that $Q_q$ is an isometry onto its image.

The finite sum in \cref{eq:BL-inverse-formula} is immediate; in the infinite
case the series converges because
$|q^{-\ell}2^{d_\ell}|_2=2^{-d_\ell}\to0$.  Writing
$z=\delta+2w$ with $\delta\in\{0,1\}$ and separating the first one when
$\delta=1$ gives the exact recursions
\begin{equation}
\label{eq:BL-Phi-recursion}
 \Phi_q(2w)=2\Phi_q(w),
 \qquad
 \Phi_q(1+2w)=\frac{2\Phi_q(w)-1}{q}.
\end{equation}
The first value is even and the second is odd.  Applying the corresponding
branch of $T_q$ therefore gives
\[
 T_q(\Phi_q(z))=\Phi_q(Sz).
\]
Iterating shows that the parity digits of $\Phi_q(z)$ are exactly the digits
of $z$, so $Q_q\Phi_q=\operatorname{id}_{\Z_2}$.  Hence $Q_q$ is
surjective; together with \cref{eq:BL-isometry-congruences}, it is bijective
and $\Phi_q=Q_q^{-1}$.  Dropping the first parity digit in
\cref{eq:BL-parity-encoder} proves $Q_qT_q=SQ_q$, and substitution of the
inverse gives the second conjugacy formula.  The congruence equivalence for
every $N$ proves that both inverse maps are solenoidal isometries and hence
homeomorphisms.
\end{proof}

Bernstein--Lagarias Appendix B is broader still.  Given arbitrary solenoidal
bijections $V_0,V_1:\Z_2\to\Z_2$, it treats
\[
 U_{V_0,V_1}(x)=
 \begin{cases}
 V_0(x/2),&x\equiv0\pmod2,\\
 V_1((x-1)/2),&x\equiv1\pmod2,
 \end{cases}
\]
and proves that its chronological parity encoder is solenoidal bijective and
$U=Q^{-1}SQ$, with a converse classification
\cite[App.~B, Thms.~B.1--B.2]{BernsteinLagarias}.  The shortened $T_q$
has $V_0(y)=y$ and $V_1(y)=qy+(q+1)/2$; the latter is an affine isometry of
$\Z_2$ for every odd $q$, with inverse
$y\mapsto q^{-1}(y-(q+1)/2)$.  This proves the exact inclusion in that
class; it does not identify Appendix B's full class with general Hydras.

\begin{theorem}[reversed purely periodic overlap]
\label{thm:BL-numen-periodic-overlap}
Let $q$ be odd, let $L\geq1$, and let
$\boldsymbol\varepsilon=(\varepsilon_0,\ldots,\varepsilon_{L-1})$ be a
chronological binary block.  If $m\geq1$, write the positions of its $m$ ones
as
\[
 0\leq v_0<\cdots<v_{m-1}<L;
\]
when $m=0$, every position list and sum below is empty.  Set
\begin{equation}
\label{eq:chronological-and-reversed-codes}
 e=\sum_{i=0}^{L-1}\varepsilon_i2^i,
 \qquad
 \operatorname{rev}_L(e)
   =\sum_{i=0}^{L-1}\varepsilon_i2^{L-1-i}.
\end{equation}
Put
\begin{equation}
\label{eq:BL-periodic-rational}
 R_{q,\boldsymbol\varepsilon}
 :=\frac{\displaystyle
         \sum_{k=0}^{m-1}q^{m-k-1}2^{v_k}}
        {2^L-q^m}\in\Q.
\end{equation}
Then
\[
 \Phi_q\!\left(\frac{e}{1-2^L}\right)
 =\iota_2\!\left(R_{q,\boldsymbol\varepsilon}\right)\in\Z_2.
\]
If, in addition, $q$ is an odd prime, so that the edition's global
$\chi_q:\Z_2\to\Z_q$ is available, then
\begin{equation}
\label{eq:BL-numen-overlap}
 \chi_q\!\left(
       \frac{\operatorname{rev}_L(e)}{1-2^L}\right)
 =\iota_q\!\left(R_{q,\boldsymbol\varepsilon}\right)\in\Z_q.
\end{equation}
Thus the two values have the same rational representative
$R_{q,\boldsymbol\varepsilon}$ after reversal.  The reversal in
\cref{eq:BL-numen-overlap} is mandatory.
\end{theorem}

\begin{proof}
If $m=0$, every displayed value is zero, so assume first that $m\geq1$.
The one positions of the purely periodic $2$-adic integer
$e/(1-2^L)$ are
\[
 d_{rm+k+1}=v_k+rL
 \qquad(r\geq0,\ 0\leq k<m).
\]
Substitution into \cref{eq:BL-inverse-formula} and a $2$-adically
convergent geometric sum give
\begin{align*}
 \Phi_q\!\left(\frac{e}{1-2^L}\right)
 &=-\sum_{r\geq0}\sum_{k=0}^{m-1}
       q^{-(rm+k+1)}2^{v_k+rL}\\
 &=-\left(\sum_{k=0}^{m-1}q^{-(k+1)}2^{v_k}\right)
       \sum_{r\geq0}\left(\frac{2^L}{q^m}\right)^r\\
 &=\iota_2\!\left(\frac{\displaystyle
         \sum_{k=0}^{m-1}q^{m-k-1}2^{v_k}}
        {2^L-q^m}\right),
\end{align*}
because $|2^L/q^m|_2=2^{-L}<1$.  This proves
the asserted $\Phi_q$ evaluation with the rational representative defined in
\cref{eq:BL-periodic-rational}.

Now put $h_0(x)=x/2$ and $h_1(x)=(qx+1)/2$.  Induction on the chronological
block gives
\begin{equation}
\label{eq:chronological-affine-return}
 h_{\varepsilon_{L-1}}\circ\cdots\circ h_{\varepsilon_0}(x)
 =\frac{q^m x+
        \displaystyle\sum_{k=0}^{m-1}q^{m-k-1}2^{v_k}}
       {2^L}.
\end{equation}
To see the induction with every index fixed, for $0\leq h\leq L$ put
\[
 m_h=\sum_{i=0}^{h-1}\varepsilon_i,
 \qquad
 C_h=\sum_{\substack{0\leq k<m\\v_k<h}}
       q^{m_h-k-1}2^{v_k}.
\]
The claim after the first $h$ chronological branches is
\[
 h_{\varepsilon_{h-1}}\circ\cdots\circ h_{\varepsilon_0}(x)
 =\frac{q^{m_h}x+C_h}{2^h}.
\]
At $h=0$ the left side denotes the empty composite
$\operatorname{id}_{\Q}$, so the claim is the identity.  For
$0\leq h<L$, assume the claim at level $h$.  If $\varepsilon_h=0$, applying
$h_0$ changes
only the denominator from $2^h$ to $2^{h+1}$, while
$m_{h+1}=m_h$ and $C_{h+1}=C_h$.  If $\varepsilon_h=1$, then
\[
 h_1\!\left(\frac{q^{m_h}x+C_h}{2^h}\right)
 =\frac{q^{m_h+1}x+qC_h+2^h}{2^{h+1}}.
\]
Here $m_{h+1}=m_h+1$; multiplication of $C_h$ by $q$ raises every old
exponent by one, and the new odd position $v_{m_h}=h$ contributes the term
$q^0 2^h$.  Thus the new numerator is exactly $q^{m_{h+1}}x+C_{h+1}$.
At $h=L$ this is \cref{eq:chronological-affine-return}.  Solving its affine
fixed-point equation gives $R_{q,\boldsymbol\varepsilon}$; its denominator is
nonzero because an odd
integer power cannot equal $2^L$.

Under Siegel's convention, the displayed word
$H_{j_0}\circ\cdots\circ H_{j_{L-1}}$ applies $j_{L-1}$ first.  It realizes
the chronological block precisely when
\[
 (j_0,\ldots,j_{L-1})
 =(\varepsilon_{L-1},\ldots,\varepsilon_0),
\]
whose little-endian code is exactly $\operatorname{rev}_L(e)$.  For prime
$q$ and $m\geq1$, its linear part is $q^m/2^L$, of $q$-adic norm $q^{-m}<1$.
The repeated-code theorem therefore evaluates its rising numen at its unique
fixed point, namely $\iota_q(R_{q,\boldsymbol\varepsilon})$.  The $m=0$ case
is again the centered
zero word.  This proves \cref{eq:BL-numen-overlap}.

The reversal is genuinely necessary in the universal formula, not merely a
different way to name the same argument.  Take $q=3$, $L=2$, and
$\boldsymbol\varepsilon=(0,1)$.  Then $e=2$ and
$\operatorname{rev}_2(e)=1$, while the formulas already proved give
\[
 \Phi_3\!\left(\frac{2}{1-4}\right)=\iota_2(2),
 \qquad
 \chi_3\!\left(\frac{1}{1-4}\right)=\iota_3(2),
 \qquad
 \chi_3\!\left(\frac{2}{1-4}\right)=\iota_3(1).
\]
Indeed, $h_1\circ h_0$ is $x\mapsto(3x+2)/4$ with fixed point $2$, whereas
$h_0\circ h_1$ is $x\mapsto(3x+1)/4$ with fixed point $1$.
Thus deleting $\operatorname{rev}_L$ makes the asserted identity false.
\end{proof}

\begin{sourceissue}[the unrestricted B\"ohm--Sontacchi Proposition 4 iterate
identity is false]
\label{status:Bohm-Sontacchi-Proposition-4}
For $q=3$, the right side of \cref{eq:chronological-affine-return} is the
formal branch-composition expression obtained by specializing
B\"ohm--Sontacchi Proposition~1.  Their printed Proposition~4 displays this
expression as $u^n(x)$ while quantifying over every $x\in\Z$ and every run
tuple $(n_0,\ldots,n_m)$.  That unrestricted left side is false.  Take
$x=0$, $m=1$, and $n_0=n_1=1$.  Then $n=3$ and $v_0=1$, so the printed right
side is $1/4$, whereas $u^3(0)=0$.  The same right side is correctly the
formal value
\[
 f\circ g\circ f(0)=\frac14,
 \qquad f(x)=\frac{x}{2},\quad g(x)=\frac{3x+1}{2}.
\]
Consequently this edition uses the formal-composition identity and identifies
it with $u^n(x)$ only after proving that the word is the actual parity
itinerary, as in \cref{prop:Tq-correct-fixed-point}.
\end{sourceissue}

The fixed point $R_{3,\boldsymbol\varepsilon}$ of that correct formal
composition is the candidate printed in B\"ohm--Sontacchi Proposition~5
\cite[p.~263]{BohmSontacchi}.  Their statement literally prints $v_m=n$;
after the explicit edition renaming $n=L$, it prints
$v_{i-1}<v_i$ for $1\leq i<m$ and $v_m=L$, but does not visibly print
$v_{m-1}<v_m$.  For $m\geq1$, the last inequality follows from the formal run
construction, where $v_m-v_{m-1}=1+n_m>0$; it is a reconstruction, not a
silent correction of Proposition~5.  The unnumbered paragraph immediately
after its proof explicitly says the candidate may lie on a cycle of length
$L$ or a proper divisor of $L$ \cite[p.~264]{BohmSontacchi}.  Thus the
independently recovered historical formula is a return-time identity, not an
exact-period classification.

\begin{corollary}[global nonidentity]
\label{cor:BL-numen-global-nonidentity}
For every odd prime $q$, the paired values with common rational representatives
in
\cref{thm:BL-numen-periodic-overlap} do not identify the global maps
$\Phi_q$ and $\chi_q$.  They have types
\[
 \Phi_q:\Z_2\to\Z_2,
 \qquad
 \chi_q:\Z_2\to\Z_q,
\]
The two values have distinct rational representatives:
\begin{equation}
\label{eq:BL-numen-value-separation}
 \Phi_q(1)=\iota_2\!\left(-\frac1q\right),
 \qquad
 \chi_q(1)=\iota_q\!\left(\frac12\right),
 \qquad
 -\frac1q\ne\frac12\quad\text{in }\Q.
\end{equation}
The map $Q_q$ must likewise retain its role as the inverse chronological
parity encoder in \cref{eq:BL-conjugacy-oriented}; it cannot be substituted
for $\Phi_q$ by a notation change.
\end{corollary}

\begin{proof}
Formula \cref{eq:BL-inverse-formula} at the single one position $d_1=0$
gives the first embedded value.  The one-letter numen is
$\chi_q(1)=H_1(0)=\iota_q(1/2)$.  Equality of the rational representatives
would force $q=-2$, contrary to oddness.  The displayed domains, codomains,
and inverse relation prove the remaining typing assertions.
\end{proof}

Laarhoven--de Weger give a graph-theoretic presentation of the same classical
conjugacy.  Their symbol $\Phi$ is the chronological graph isomorphism
satisfying $T=\Phi^{-1}\sigma\Phi$, so in the notation fixed here it plays
the role of $Q_3$, not $\Phi_3$.  Their odd-$an+b$ analogue is exact, whereas
their $p$-ary discussion uses ``appropriately chosen'' coefficients and does
not cover arbitrary Hydras
\cite[Thms.~3.1 and~5.1 and Sec.~5.2]{LaarhovenDeWeger}.

\subsection{Branch correctness for \texorpdfstring{$T_q$}{Tq}}

The shortened map admits a particularly clean correctness argument in
$\Q_2$.  Extend the two affine formulas in \cref{eq:Tq-branches} to all of
$\Q_2$.

\begin{lemma}[one correct branch]
\label{lem:one-correct-Tq-branch}
Let $q$ be odd and $x\in\Z_2$.  Exactly one of $H_0(x),H_1(x)$ belongs to
$\Z_2$, namely $H_{[x]_2}(x)$.  If a wrong branch is applied, the result has
$2$-adic valuation $-1$.  If $v_2(y)<0$, then
\[
 v_2(H_0(y))=v_2(H_1(y))=v_2(y)-1.
\]
\end{lemma}

\begin{proof}
If $x$ is even, $x/2\in\Z_2$ and $qx+1$ is odd, so
$(qx+1)/2\notin\Z_2$.  If $x$ is odd, $x/2\notin\Z_2$ and $qx+1$ is even,
so $(qx+1)/2\in\Z_2$.  This proves the first two assertions.  If
$v_2(y)<0$, then $v_2(qy)=v_2(y)<0=v_2(1)$; the ultrametric inequality with
unequal valuations gives $v_2(qy+1)=v_2(y)$.  Division by $2$ lowers either
valuation by one.
\end{proof}

\begin{proposition}[formal fixed points in $\Z_2$ are genuine]
\label{prop:Tq-correct-fixed-point}
Let $\mathbf j$ be a finite binary word.  If $z\in\Z_2$ and
$H_{\mathbf j}(z)=z$, then every branch in the right-to-left evaluation of
$H_{\mathbf j}(z)$ is correct.  Consequently
\[
 T_q^{\circ|\mathbf j|}(z)=z.
\]
\end{proposition}

\begin{proof}
Start with $z\in\Z_2$ and evaluate the rightmost branch first.  If any branch
is wrong, \cref{lem:one-correct-Tq-branch} says the intermediate value leaves
$\Z_2$ with negative valuation; every later branch decreases that valuation
further.  The final value then cannot equal $z\in\Z_2$.  Hence no branch was
wrong, and the formal composition is the genuine iterate with the reversed
chronological itinerary, as in \cref{eq:chronological-composition}.
\end{proof}

\subsection{The checked \texorpdfstring{$T_q$}{Tq}
periodic-point correspondence}

Assume now that $q$ is an odd prime, so that the numen
$\chi_q:\Z_2\to\Z_q$ is defined everywhere by
\cref{eq:chi-q-adic-series}.

\begin{theorem}[Siegel's periodic-word correspondence for $T_q$]
\label{thm:Tq-periodic-word-CP}
Let $n\geq1$, let $L=\lambda_2(n)$, and let $\mathbf j$ be the shortest
least-significant-digit-first binary word for $n$.  Then
\begin{equation}
\label{eq:chi-B2}
 \chi_q(B_2(n))
 =\frac{\chi_q(n)}{1-M_q(n)},
 \qquad
 M_q(n)=\frac{q^{\#_1(n)}}{2^L}.
\end{equation}
The right side is a rational number.  If it belongs to $\Z_2$, it is a
periodic point of the $2$-adic extension of $T_q$; if it belongs to $\Z$, it
is a periodic integer point of $T_q$.

Conversely, let $x$ be a nonzero periodic integer point of exact period $L$.
If $\mathbf j$ is the length-$L$ word obtained by reversing one complete
chronological itinerary of $x$ and $n=\Dig(\mathbf j)$, then
\begin{equation}
\label{eq:padded-cycle-code}
 x=\chi_q\bigl(B_{2,L}(n)\bigr),
 \qquad B_{2,L}(n)=\frac{n}{1-2^L}.
\end{equation}
At least one cyclic phase of the cycle has a reversed word whose final digit
is $1$; for that phase $L=\lambda_2(n)$ and
$B_{2,L}(n)=B_2(n)$.  Other phases may require the padded length-$L$ code.
\end{theorem}

\begin{proof}
The word contains at least one $1$, hence
$|M_q(n)|_q=q^{-\#_1(n)}<1$.  Apply
\cref{thm:repeated-code-fixed-point}; in the scalar case it gives exactly
\cref{eq:chi-B2}.  If the fixed point lies in $\Z_2$,
\cref{prop:Tq-correct-fixed-point} turns the formal fixed-point equation into
a genuine periodicity equation.

Conversely, let $x\ne0$ have period $L$ and let
$(\varepsilon_0,\ldots,\varepsilon_{L-1})$ be its chronological parity
itinerary.  Put
$\mathbf j=(\varepsilon_{L-1},\ldots,\varepsilon_0)$, so
$H_{\mathbf j}(x)=T_q^{\circ L}(x)=x$.  The word contains a $1$, since an
all-zero cycle would satisfy $x/2^L=x$ and hence $x=0$.  The affine fixed
point is unique because
$M_q(\mathbf j)=q^{\#_1(\mathbf j)}/2^L\ne1$: an odd power of the odd
integer $q$ cannot equal a positive power of $2$.  Therefore
$x=(1-M_q(\mathbf j))^{-1}\chi_q(\mathbf j)$, which is
$\chi_q(B_{2,L}(n))$ by \cref{thm:repeated-code-fixed-point}.  Rotating the
chronological starting point rotates the word.  One can choose a rotation
whose high digit is nonzero and thereby use the shortest-word abbreviation
$B_2(n)$, but that rotation generally changes $x$ to another point of its
cycle.
\end{proof}

\begin{sourceissue}[shortest words do not encode every cycle phase]
The stronger version in the 2007 source overclaims that every point in the
numen image of a cycle is represented by $B_2(n)$ for the shortest positive
word of $n$.  For $T_3$, the point $2$ on the cycle $1\leftrightarrow2$ is
encoded by the length-two word $(1,0)$; its high zero is part of the repeated
block, even though the ordinary integer represented by that word is $1$.
The synchronized code is $B_{2,2}(1)=1/(1-4)=-1/3$, not
$B_2(1)=1/(1-2)=-1$.  The padded formula
\cref{eq:padded-cycle-code} preserves the phase.  The source also writes a
truncation modulus $2^m$ for $m$ repeated length-$L$ blocks where the typed
modulus is $2^{mL}$; taking the correct subsequence repairs the limit.
\end{sourceissue}

\begin{sourceissue}[the centered zero orbit]
The 2026 Correspondence Principle is displayed for every periodic
$z\in\OK$ but requires a preimage in
$(\Z_p\setminus\N_0)\cap\Q$.  A centered Hydra has the periodic point $0$,
encoded by the all-zero word, whose digit-ring element is $0\in\N_0$.
No general argument in the source produces a non-eventually-zero rational
code with numen value $0$.  The earlier $T_q$ theorem correctly treats
nonzero periodic points.  The general statement therefore needs either an
explicit exception for $0$, inclusion of the code $0$, or an additional
surjectivity assertion about the zero fibre.
\end{sourceissue}

\subsection{The general 2026 Correspondence Principle}

We now record the latest source theorem without upgrading its proof status.
Let $H:\OK\to\OK$ be a proper, centered, exact-integral $\Lambda$-Hydra and
$p=|\OK/\Lambda|$.  Normalize a finite-place additive valuation by
$v(K^\times)=\Z$ and write
\[
 v(\Lambda)=\min\{v(\lambda):0\ne\lambda\in\Lambda\}.
\]
The source assumes:
\begin{enumerate}[label=\textup{(H\arabic*)}]
\item for each integer $n\geq1$ there is a place $\ell_n$ of $K$ such that
      $\|M_H(n)\|_{\ell_n}<1$;
\item there is a nontrivial nonarchimedean place $v$ such that
      $\|r_j\|_v>1$ for every branch and $v(\Lambda)\geq1$.
\end{enumerate}

\begin{status}[General CP (source)]
Under (H1)--(H2), arXiv:2601.17030v3 asserts:
\begin{enumerate}[label=\textup{(\roman*)}]
\item an integral point is periodic if and only if it is the value of $X_H$
      at a rational, non-eventually-zero base-$p$ digit input;
\item if $K$ is a number field and an irrational digit input has numen value
      $x\in\OK$, then
      $|H^{\circ m}(x)|_w\to\infty$ as $m\to\infty$ for every nontrivial
      archimedean absolute value $|\cdot|_w$.
\end{enumerate}
The source labels its proof ``Sketch.''  The repeated-word affine calculation
needed for the forward direction is proved in full in
\cref{thm:repeated-code-fixed-point}.  The converse still depends on the
source's completion-level ``correctness'' argument, and the divergence clause
depends on an identification of fibres not established by the displayed
hypotheses.  Accordingly, this edition records (i)--(ii) as source claims,
not as reconstructed theorems.
\end{status}

\begin{sourceissue}[defects inside the 2026 proof sketch]
The sketch contains five points which matter mathematically.
\begin{enumerate}[label=\textup{(\alph*)}]
\item It writes $M_q$ in a theorem about a general $H$; the typed symbol is
      $M_H$.
\item It writes a quotient $(1-M_H^m)/(1-M_H)$ although $M_H$ is generally
      an endomorphism.  The typed expression is
      $(I-M_H^m)(I-M_H)^{-1}$, and invertibility must be proved.
\item The proof of the divergence clause invokes a hypothesis ``(iii)'', but
      the theorem lists only (i) and (ii).
\item From the periodic-point correspondence, existence of a rational code
      for a periodic value does not imply that every code in the same numen
      fibre is rational.  Such an inference requires injectivity or a precise
      fibre theorem, neither of which is part of the displayed statement.
\item In dimension greater than one, $\|r_j\|_v>1$ is only a statement about
      the \emph{largest} expansion of $r_j$.  It does not imply
      $|r_jx|_v>|x|_v$ for every nonzero $x$, which is the kind of
      branch-detection estimate used by the sketch.  A sufficient typed
      replacement is a lower expansion bound
      $\inf_{x\ne0}|r_jx|_v/|x|_v>1$, equivalently
      $\|r_j^{-1}\|_v<1$ for invertible $r_j$, together with the required
      translation/integrality control.  In one scalar dimension the upper
      and lower factors coincide, but the general endomorphism statement
      cannot use the operator supremum alone.
\end{enumerate}
These are not stylistic complaints; they delimit what has and has not been
proved by the sketch.
\end{sourceissue}

\subsection{Weak Collatz versus full Collatz}

For $T_3$, the \emph{weak Collatz conjecture} in Siegel's contour paper is a
periodic-orbit assertion: the only integer cycles are the known ones (with
the sign and state-space convention stated in that paper).  The full Collatz
conjecture additionally requires every positive integer orbit to enter the
positive $1\leftrightarrow2$ cycle.  A characterization of periodic points,
even a complete one, does not rule out divergent positive trajectories.
The contour equivalence in \cref{sec:contour-mellin} is therefore kept under
the name \emph{weak}; it is not promoted to an equivalence with the full
conjecture.
